%% toc: P1 7 7.1,7.2,7.3
\section{Ambiguity, restricted quantification, uniqueness, and scope in translation}
\label{sec:p01-w02-p01-022-translation}

The method of \cref{sec:levels-05} tells the translator to resolve pronouns,
plurality, restricted quantification and uniqueness before translating, and
to keep quantifier scope with parentheses
\cite[ch.~``Levels, Objects, and Consequence'', sec.~``Arguments and
formalization'', method ``Formalizing an argument'']{FrozenFormalLogicFormalizing}.
It does not say how an ambiguity is found, how a reading is chosen, how the
choice is written down, which formula a restricted quantifier phrase or a
uniqueness phrase becomes, or what goes wrong when the wrong pattern is used.
This section supplies those parts of the translation step.

\emph{Inputs.} A natural-language argument, a declared domain, a declared
first-order signature \(\Lang\) with logical equality
(\cref{def:p02-fol-signature}), and a declared set of background sentences.
The semantics is the satisfaction relation of \cref{sec:fol-semantics-02};
substitution \(\varphi[t/x]\) is the capture-avoiding substitution of
\cref{sec:fol-syntax-semantics-04}; consequence \(\Gamma\entails\varphi\) and
logical equivalence of formulas are those of the definition ``Truth,
validity, and semantic consequence'' in \cref{sec:fol-semantics-02}
\cite[ch.~``Structures, Satisfaction, and Translation'', secs.~``Satisfaction''
and ``Translation with quantifiers'']{FrozenFormalLogicFormalizing}.

\emph{Conclusions.} A record of readings and a test that separates harmless
from harmful ambiguity (\cref{def:p01-w02-p01-022-reading,thm:p01-w02-p01-022-readings,def:p01-w02-p01-022-ambiguity-method}); the
semantics of restricted quantifiers and the exact cost of the two standard
mistranslations (\cref{thm:p01-w02-p01-022-relativization,thm:p01-w02-p01-022-mismatch}); exact formulas for ``at most one'',
``exactly one'' and ``the'', with their counting characterization and the
freshness condition they need (\cref{def:p01-w02-p01-022-uniqueness,thm:p01-w02-p01-022-uniqueness,thm:p01-w02-p01-022-freshness}); and the
scope behaviour of eliminated definite descriptions
(\cref{thm:p01-w02-p01-022-description-scope}). The results and their proofs
are editorial completions built on the inherited definitions. The contextual
analysis of descriptions and the conditional reading of ``all'' are
Russell's \cite[481--482, 489--490]{Russell1905}; the proofs given here are
not attributed to him.

\subsection{Setting}
\label{sec:p01-w02-p01-022-setting}

\begin{hbconvention}{def:p01-w02-p01-022-setting}{Background sets and solution sets}
Throughout, \(\Lang\) is a declared signature, \(\mathcal M\) ranges over
\(\Lang\)-structures with nonempty domain \(\dom\mathcal M\), and \(s\)
ranges over variable assignments in \(\mathcal M\), with update
\(s[x\mapsto a]\).
\begin{enumerate}
  \item A \emph{background set} is a set \(\Pi\) of \(\Lang\)-sentences fixed
        before translation. It records the connections between declared
        symbols that the intended reading takes for granted, for example
        that no paper is a referee. The argument
        \(\langle\Gamma,\psi\rangle\) is \emph{valid relative to \(\Pi\)}
        iff \(\Gamma\cup\Pi\entails\psi\). With \(\Pi=\varnothing\) this is
        semantic validity in the sense of \cref{sec:levels-05}.
  \item Two \(\Lang\)-formulas \(\varphi_1,\varphi_2\) are
        \emph{\(\Pi\)-equivalent} iff
        \(\mathcal M,s\entails\varphi_1\Leftrightarrow
        \mathcal M,s\entails\varphi_2\) for every model \(\mathcal M\) of
        \(\Pi\) and every assignment \(s\). For \(\Pi=\varnothing\) this is
        logical equivalence.
  \item For a formula \(\varphi\), a variable \(x\), a structure
        \(\mathcal M\) and an assignment \(s\), the \emph{solution set} of
        \(\varphi\) in \(x\) is
        \[
          \varphi^{\mathcal M,s}_x\defeq
          \{a\in\dom\mathcal M:\mathcal M,s[x\mapsto a]\entails\varphi\}.
        \]
        By formula coincidence (\cref{lem:fol-coincidence}) it depends on
        \(s\) only through the values of \(s\) on
        \(\FV(\varphi)\setminus\{x\}\).
\end{enumerate}
\end{hbconvention}

\subsection{Readings and ambiguity}
\label{sec:p01-w02-p01-022-ambiguity}

\begin{hbdefinition}{def:p01-w02-p01-022-reading}{Dictionary, reading, ambiguity}
Fix a domain, a signature \(\Lang\) and a background set \(\Pi\).
\begin{enumerate}
  \item A \emph{dictionary} assigns to each content word or phrase of the
        argument a declared symbol of \(\Lang\) together with the sense in
        which the symbol is meant.
  \item A \emph{reading} of a natural-language sentence is an
        \(\Lang\)-formula together with a statement, in the metalanguage, of
        which interpretation of the sentence it formalizes. A reading of a
        premise or conclusion is a sentence of \(\Lang\) unless a free
        variable is declared as a parameter.
  \item A sentence is \emph{ambiguous relative to \(\Lang\) and \(\Pi\)}
        iff it has two readings that are not \(\Pi\)-equivalent. Two such
        readings are \emph{separated} by any model \(\mathcal M\) of \(\Pi\)
        and assignment \(s\) at which exactly one of them is satisfied; this
        pair is a \emph{separating structure}.
  \item An ambiguity is \emph{harmless for an argument} iff the verdict
        ``valid relative to \(\Pi\)'' or ``invalid relative to \(\Pi\)'' is
        the same under every choice of readings; otherwise it is
        \emph{harmful}.
\end{enumerate}
\index{reading}\index{ambiguity!relative to a signature}
\end{hbdefinition}

Ambiguity in this sense is a relation between a sentence and a declared
language. A sentence whose two interpretations become \(\Pi\)-equivalent
formulas is not ambiguous for the argument at hand, however different the
interpretations sound; and a sentence that reads unambiguously in ordinary
speech may still have several readings once the signature forces a choice,
for example between a one-place predicate of individuals and a relation to
a collection.

\begin{hbdefinition}{def:p01-w02-p01-022-kinds}{Kinds of ambiguity and their formal effect}
The following kinds are distinguished by what changes between the readings.
\begin{enumerate}
  \item \emph{Lexical.} One word has two senses, and the dictionary offers
        two symbols. The readings differ in a nonlogical symbol.
  \item \emph{Structural.} The sentence has two grammatical analyses, for
        example two groupings of ``and'' and ``or'' or two attachments of a
        modifying clause. The readings differ in the formation tree of the
        formula.
  \item \emph{Scope.} The sentence contains two or more scope-bearing items
        (quantifier phrases, negation, ``only'', definite descriptions) whose
        relative order is not fixed by the wording. The readings differ in
        the order in which the corresponding operators occur on a branch of
        the formation tree.
  \item \emph{Pronoun.} A pronoun has more than one candidate antecedent, or
        may either refer to a fixed object or be bound by a quantifier
        phrase. The readings differ in which variable or constant occupies
        the pronoun's position, and in which quantifier binds that variable.
  \item \emph{Plurality.} A plural noun phrase combines with a predicate
        either distributively (the predicate holds of each member) or
        collectively (the predicate holds of the members taken together), or
        its grammatical number may or may not be meant to assert more than
        one object. The readings differ in whether the predicate applies to
        individuals or to a collection, and in whether a distinctness
        condition is asserted.
\end{enumerate}
\index{ambiguity!lexical}\index{ambiguity!structural}\index{ambiguity!scope}
\index{ambiguity!pronoun}\index{ambiguity!plurality}
\end{hbdefinition}

\begin{hbtheorem}{thm:p01-w02-p01-022-readings}{proved=annex:p01-w02-p01-022-readings}{Validity under alternative readings}
Let \(\Pi\) be a background set, \(\Gamma\) a set of \(\Lang\)-formulas,
\(\varphi_1,\ldots,\varphi_k\) (\(k\geq1\)) readings of one further premise
and \(\psi_1,\ldots,\psi_m\) (\(m\geq1\)) readings of the conclusion.
\begin{enumerate}
  \item If \(\varphi_1\) and \(\varphi_2\) are \(\Pi\)-equivalent, then for
        every formula \(\psi\),
        \(\Gamma\cup\Pi\cup\{\varphi_1\}\entails\psi\) iff
        \(\Gamma\cup\Pi\cup\{\varphi_2\}\entails\psi\). If \(\psi_1\) and
        \(\psi_2\) are \(\Pi\)-equivalent, then
        \(\Gamma\cup\Pi\entails\psi_1\) iff \(\Gamma\cup\Pi\entails\psi_2\).
  \item For every formula \(\psi\),
        \(\Gamma\cup\Pi\cup\{\varphi_i\}\entails\psi\) for every
        \(i\leq k\) iff
        \(\Gamma\cup\Pi\cup\{\varphi_1\lor\cdots\lor\varphi_k\}\entails\psi\).
  \item \(\Gamma\cup\Pi\entails\psi_j\) for every \(j\leq m\) iff
        \(\Gamma\cup\Pi\entails\psi_1\land\cdots\land\psi_m\).
\end{enumerate}
\end{hbtheorem}

Clause 1 says that only non-equivalent readings need to be recorded
separately. Clauses 2 and 3 say when an unresolved ambiguity is harmless for
a verdict of validity: the argument is valid under every combination of
readings iff it is valid with the disjunction of the premise readings and the
conjunction of the conclusion readings. When the argument is valid under some
combinations and invalid under others, the ambiguity is harmful and the
verdict has to be reported reading by reading.

\begin{hbmethod}{def:p01-w02-p01-022-ambiguity-method}{Detecting, resolving and recording ambiguity}
Apply the following steps to each premise and to the conclusion after the
domain, the signature and \(\Pi\) have been declared and before the formula
is written.
\begin{enumerate}
  \item \emph{Lexical check.} For each content word, list the dictionary
        entries it could take. More than one candidate symbol is a lexical
        ambiguity point.
  \item \emph{Structural check.} Group the clauses explicitly with
        parentheses. If two groupings, or two attachments of a modifying
        clause, are grammatical, record a structural ambiguity point.
  \item \emph{Scope check.} List the scope-bearing items of each clause.
        For each pair whose order is not fixed by the wording, write both
        orders. Each distinct order is a candidate reading.
  \item \emph{Pronoun check.} For each pronoun, list the candidate
        antecedents and decide, for each, whether the pronoun names a fixed
        object (a constant) or is bound by a quantifier phrase (a variable).
        A pronoun whose candidate binder does not have it in its scope cannot
        be translated by that binder as written
        (\cref{thm:p01-w02-p01-022-indefinite}).
  \item \emph{Plurality check.} For each plural noun phrase, ask whether its
        predicate is meant of each member or of the members together, and
        whether the plural is meant to assert more than one object. A
        collective reading needs a declared domain of collections and a
        membership relation; a ``more than one'' reading needs a
        distinctness condition.
  \item \emph{Separation.} Translate each candidate reading. For each pair
        of readings, either exhibit a separating model of \(\Pi\) and
        assignment, or show that the pair is \(\Pi\)-equivalent. Equivalent
        readings are merged by clause 1 of \cref{thm:p01-w02-p01-022-readings}.
  \item \emph{Resolution.} Choose among the remaining readings using
        evidence that is stated: the other premises, the declared purpose of
        the argument, and what its author says elsewhere. A reading under
        which a premise is false in every intended model is evidence against
        that reading, not a reason to reject the argument. The same word is
        read in the same sense in every premise unless a change of sense is
        recorded.
  \item \emph{Recording.} Write a translation record containing the domain,
        the signature, \(\Pi\), the dictionary, and, for each ambiguity point,
        its kind, the readings, the separating structures, the chosen reading
        and the evidence for the choice.
  \item \emph{Unresolved points.} If a point cannot be resolved, determine
        validity for each combination of readings, using clauses 2 and 3 of
        \cref{thm:p01-w02-p01-022-readings} to test all combinations at once,
        and report the verdict per reading.
\end{enumerate}
\index{ambiguity!method}\index{translation record}
\end{hbmethod}

\begin{hbexample}{def:p01-w02-p01-022-lexical}{Lexical ambiguity and equivocation}
Consider: every bank is beside a river; Halden is a bank; so Halden is
beside a river. The dictionary offers two symbols for ``bank'':
\(B_1(x)\), ``\(x\) is a riverbank'', and \(B_2(x)\), ``\(x\) is a
financial institution''. Let \(W(x)\) mean ``\(x\) is beside a river'' and
let \(h\) be a constant for Halden. The uniform readings
\[
  \forall x\,(B_i(x)\to W(x)),\ B_i(h)\ \entails\ W(h)\qquad(i=1,2)
\]
are valid: in any structure satisfying both premises, \(h^{\mathcal M}\) is
in \(B_i^{\mathcal M}\) and so in \(W^{\mathcal M}\). The mixed reading
\(\forall x\,(B_1(x)\to W(x)),\ B_2(h)\) with conclusion \(W(h)\) is invalid:
take \(\dom\mathcal M=\{0\}\), \(h^{\mathcal M}=0\),
\(B_1^{\mathcal M}=W^{\mathcal M}=\varnothing\) and
\(B_2^{\mathcal M}=\{0\}\). The first premise holds vacuously, the second
holds, and the conclusion fails. The argument is persuasive only under the
mixed reading, where the first premise is plausible in the riverbank sense
and the second in the financial sense; this is equivocation, and step 7 of
\cref{def:p01-w02-p01-022-ambiguity-method} forbids it unless the change of
sense is recorded, in which case the record exhibits the invalidity.
\end{hbexample}

\begin{hbtheorem}{thm:p01-w02-p01-022-negation-scope}{proved=annex:p01-w02-p01-022-negation-scope}{Negation scope over a restricted universal}
Let \(S\) and \(P\) be unary relation symbols and put
\[
  \rho_1\defeq\forall x\,(S(x)\to\neg P(x)),\qquad
  \rho_2\defeq\neg\forall x\,(S(x)\to P(x)).
\]
\begin{enumerate}
  \item \(\rho_1\not\entails\rho_2\) and \(\rho_2\not\entails\rho_1\).
  \item \(\rho_1,\exists x\,S(x)\entails\rho_2\).
\end{enumerate}
\end{hbtheorem}

\begin{hbexample}{def:p01-w02-p01-022-scope-example}{A scope ambiguity and its record}
``Every student did not pass'', with \(S\) for student and \(P\) for passed,
has the two readings \(\rho_1\) (no student passed) and \(\rho_2\) (not every
student passed) of \cref{thm:p01-w02-p01-022-negation-scope}. The structure
with domain \(\{0,1\}\), \(S^{\mathcal M}=\{0,1\}\) and
\(P^{\mathcal M}=\{0\}\) separates them, so they must be recorded
separately. If the context says that some students passed, \(\rho_1\)
contradicts the context, which is evidence for \(\rho_2\). If the argument
needs only the weaker \(\rho_2\) and the premises guarantee that there is a
student, clause 2 of the theorem shows that the choice is harmless for that
inference.
\end{hbexample}

\begin{hblemma}{thm:p01-w02-p01-022-indefinite}{proved=annex:p01-w02-p01-022-indefinite}{Indefinites in conditional antecedents}
Let \(\varphi\) and \(\chi\) be formulas and \(y\) a variable with
\(y\notin\FV(\chi)\). Then \(\exists y\,\varphi\to\chi\) and
\(\forall y\,(\varphi\to\chi)\) are logically equivalent.
\end{hblemma}

\begin{hbexample}{def:p01-w02-p01-022-pronoun}{A pronoun bound by an indefinite}
Take ``every referee who read a paper graded it'', with \(R(x)\) referee,
\(D(y)\) paper, \(H(x,y)\) ``\(x\) read \(y\)'' and \(G(x,y)\) ``\(x\) graded
\(y\)''. The literal rendering
\[
  \forall x\,\bigl((R(x)\land\exists y\,(D(y)\land H(x,y)))\to G(x,y)\bigr)
\]
is not a sentence: the occurrence of \(y\) in \(G(x,y)\) lies outside the
scope of \(\exists y\), so the pronoun is not bound by its antecedent.
Widening the existential to
\(\forall x\,\exists y\,((R(x)\land D(y)\land H(x,y))\to G(x,y))\) binds the
pronoun but produces the mismatch pattern of
\cref{thm:p01-w02-p01-022-mismatch}: in a model of a background set containing
\(\forall x\,\neg(R(x)\land D(x))\), for every \(x\) the choice
\(y\mapsto x\) makes the antecedent false, because nothing is both a
referee and a paper; so the sentence is true whatever anyone graded. The
recorded reading is
\[
  \forall x\,\forall y\,\bigl((R(x)\land D(y)\land H(x,y))\to G(x,y)\bigr).
\]
\Cref{thm:p01-w02-p01-022-indefinite} explains the universal force: when the
consequent does not mention the paper, the indefinite in the antecedent is
equivalent to a universal over the whole conditional, and the recorded
reading is the extension of that universal form in which the same variable
also fills the pronoun's position. This is a translation choice, stated in
the record; it is not derived from the literal rendering, which is not a
sentence.
\end{hbexample}

\begin{hbexample}{def:p01-w02-p01-022-plurality}{Plural readings}
``The students lifted the piano.'' With \(S(x)\) student, \(L(x,y)\) ``\(x\)
lifted \(y\)'' and a constant \(p\) for the piano, the distributive reading
is \(\forall x\,(S(x)\to L(x,p))\), conjoined with \(\exists x\,S(x)\) if the
plural is taken to assert that there are students. A collective reading
cannot be stated with \(L\) applied to individuals; declare a domain that
also contains collections, a relation \(\mathit{In}(x,g)\) for ``\(x\) is a
member of \(g\)'', and read
\(\exists g\,\bigl(\forall x\,(\mathit{In}(x,g)\leftrightarrow S(x))\land
L(g,p)\bigr)\). The structure with domain \(\{u,v,g,p\}\), students
\(u,v\), \(\mathit{In}^{\mathcal M}=\{(u,g),(v,g)\}\) and
\(L^{\mathcal M}=\{(g,p)\}\) satisfies the collective reading and not the
distributive one; replacing \(L^{\mathcal M}\) by \(\{(u,p),(v,p)\}\)
reverses this. Number marking is a separate question. Russell set aside the
suggestion of ``only one'' in ``a man'' and of ``more than one'' in ``some
men'' when analysing these phrases \cite[481, footnote]{Russell1905};
\(\exists x\,(S(x)\land L(x,p))\) asserts at least one. If ``some students''
is meant to assert at least two, the reading is
\(\exists x\,\exists y\,(x\ne y\land S(x)\land S(y)\land L(x,p)\land
L(y,p))\).
\end{hbexample}

\subsection{Restricted quantification}
\label{sec:p01-w02-p01-022-restricted}

\Cref{sec:fol-semantics-03} introduces \(\forall x\in A\,\varphi\) and
\(\exists x\in A\,\varphi\) as abbreviations of
\(\forall x\,(A(x)\to\varphi)\) and \(\exists x\,(A(x)\land\varphi)\)
\cite[ch.~``Structures, Satisfaction, and Translation'', sec.~``Translation
with quantifiers'']{FrozenFormalLogicFormalizing}. Russell gives the same
analysis of ``all men are mortal'' as a hypothetical that holds of every
object, and of ``a man'' as a conjunction that is not always false
\cite[481]{Russell1905}. Below the restrictor may be any formula
\(\alpha\), so that ``every referee who read a paper'' is restricted by a
compound condition.

\begin{hbtheorem}{thm:p01-w02-p01-022-relativization}{proved=annex:p01-w02-p01-022-relativization}{Semantics of restricted quantifiers}
Let \(\alpha\) and \(\varphi\) be formulas, \(x\) a variable, \(\mathcal M\) a
structure and \(s\) an assignment, and write
\(A=\alpha^{\mathcal M,s}_x\).
\begin{enumerate}
  \item \(\mathcal M,s\entails\forall x\,(\alpha\to\varphi)\) iff
        \(\mathcal M,s[x\mapsto a]\entails\varphi\) for every \(a\in A\).
  \item \(\mathcal M,s\entails\exists x\,(\alpha\land\varphi)\) iff
        \(\mathcal M,s[x\mapsto a]\entails\varphi\) for some \(a\in A\).
  \item \(\neg\forall x\,(\alpha\to\varphi)\) is logically equivalent to
        \(\exists x\,(\alpha\land\neg\varphi)\), and
        \(\neg\exists x\,(\alpha\land\varphi)\) is logically equivalent to
        \(\forall x\,(\alpha\to\neg\varphi)\).
  \item If \(A=\varnothing\), then \(\mathcal M,s\entails\forall x\,
        (\alpha\to\varphi)\) and \(\mathcal M,s\not\entails\exists x\,
        (\alpha\land\varphi)\).
\end{enumerate}
\end{hbtheorem}

Clauses 1 and 2 say that the conditional and the conjunction are exactly
what make the unrestricted quantifier range over \(A\) alone; the
connective is not a stylistic choice. Clause 3 is quantifier duality inside
the restriction: the negation of a restricted universal is a restricted
existential, with the restrictor unchanged.

\begin{hbmethod}{def:p01-w02-p01-022-patterns}{Translating restricted forms}
With \(\alpha\) the formula for the restricting noun phrase and \(\varphi\)
the formula for the predicate, both with \(x\) free, translate as follows.
\begin{center}
\begin{tabular}{@{}ll@{}}
\toprule
Form & Formula \\
\midrule
every \(\alpha\) is \(\varphi\) & \(\forall x\,(\alpha\to\varphi)\) \\
some \(\alpha\) is \(\varphi\) & \(\exists x\,(\alpha\land\varphi)\) \\
no \(\alpha\) is \(\varphi\) & \(\forall x\,(\alpha\to\neg\varphi)\),
  equivalently \(\neg\exists x\,(\alpha\land\varphi)\) \\
not every \(\alpha\) is \(\varphi\) & \(\exists x\,(\alpha\land\neg\varphi)\),
  equivalently \(\neg\forall x\,(\alpha\to\varphi)\) \\
only \(\alpha\) are \(\varphi\) & \(\forall x\,(\varphi\to\alpha)\) \\
every \(\alpha\) that is \(\beta\) is \(\varphi\) &
  \(\forall x\,((\alpha\land\beta)\to\varphi)\) \\
\bottomrule
\end{tabular}
\end{center}
The equivalences in the third and fourth rows are clause 3 of
\cref{thm:p01-w02-p01-022-relativization}. In ``only \(\alpha\) are
\(\varphi\)'' the restricting phrase supplies the consequent, as ``only if''
does in the propositional case of \cref{sec:prop-syntax-semantics-03}. If
``every \(\alpha\) is \(\varphi\)'' is meant to imply that there are
\(\alpha\)s, add \(\exists x\,\alpha\) as a separate conjunct and record the
addition; the universal does not carry it
(\cref{thm:p01-w02-p01-022-mismatch}, clauses 5 and 6).
\index{restricted quantification!translation patterns}
\end{hbmethod}

\begin{hbtheorem}{thm:p01-w02-p01-022-mismatch}{proved=annex:p01-w02-p01-022-mismatch}{The conditional and conjunction mismatches}
Let \(\alpha,\varphi\) be formulas and \(x\) a variable.
\begin{enumerate}
  \item \(\forall x\,(\alpha\land\varphi)\) is logically equivalent to
        \(\forall x\,\alpha\land\forall x\,(\alpha\to\varphi)\).
  \item \(\exists x\,(\alpha\to\varphi)\) is logically equivalent to
        \(\exists x\,\neg\alpha\lor\exists x\,(\alpha\land\varphi)\).
\end{enumerate}
For unary relation symbols \(A\) and \(Q\):
\begin{enumerate}\setcounter{enumi}{2}
  \item \(\forall x\,(A(x)\land Q(x))\entails\forall x\,(A(x)\to Q(x))\),
        and the converse fails.
  \item \(\exists x\,(A(x)\land Q(x))\entails\exists x\,(A(x)\to Q(x))\),
        and the converse fails.
  \item \(\forall x\,(A(x)\to Q(x))\not\entails\exists x\,(A(x)\land Q(x))\).
  \item \(\forall x\,(A(x)\to Q(x)),\ \exists x\,A(x)\entails
        \exists x\,(A(x)\land Q(x))\).
\end{enumerate}
\end{hbtheorem}

\begin{hbwarning}{def:p01-w02-p01-022-mismatch-warning}{Conditional with the existential, conjunction with the universal}
Writing ``every \(A\) is \(Q\)'' as \(\forall x\,(A(x)\land Q(x))\) asserts in
addition that everything in the domain is an \(A\)
(\cref{thm:p01-w02-p01-022-mismatch}, clause 1). Over a domain that contains
anything other than \(A\)s, for instance referees and papers together, the
sentence is false whatever \(Q\) is. Writing ``some \(A\) is \(Q\)'' as
\(\exists x\,(A(x)\to Q(x))\) is satisfied by any object that is not an
\(A\) (clause 2), so over such a domain it is true whatever \(Q\) is. Both
errors make the translation insensitive to the predicate, which is the test
to apply: a correct restricted translation changes truth value when the
predicate's extension on the restrictor changes and is unaffected by objects
outside the restrictor (\cref{thm:p01-w02-p01-022-relativization}, clauses 1
and 2). The error taxonomy of \cref{sec:writing-foundations-07} lists
explicit restricted quantification as the repair for a scope error; this
warning concerns the connective inside the restriction.
\index{restricted quantification!mismatch}
\end{hbwarning}

A restricted quantifier relativizes the domain without a change of
language. Strict many-sorted logic does the same with a sort per restrictor
and reduces to one-sorted logic by relativizing quantifiers to sort
predicates (\cref{sec:p02-many-sorted}).

\subsection{Uniqueness and counting}
\label{sec:p01-w02-p01-022-uniqueness}

\begin{hbdefinition}{def:p01-w02-p01-022-uniqueness}{At most one, exactly one}
Let \(\varphi\) be a formula and \(x\) a variable. A \emph{comparison
variable} for \((\varphi,x)\) is a variable \(y\ne x\) that does not occur in
\(\varphi\), free or bound. For a comparison variable \(y\), define
\[
\begin{aligned}
  \exists^{\leq1}x\,\varphi
  &\defeq\forall x\,\forall y\,\bigl((\varphi\land\varphi[y/x])\to x=y\bigr),\\
  \exists!x\,\varphi
  &\defeq\exists x\,\bigl(\varphi\land\forall y\,(\varphi[y/x]\to y=x)\bigr).
\end{aligned}
\]
They are read ``at most one \(x\) is such that \(\varphi\)'' and ``exactly one
\(x\) is such that \(\varphi\)''. Different comparison variables give
logically equivalent formulas (\cref{thm:p01-w02-p01-022-uniqueness}), so the
choice is not recorded. These are metalinguistic abbreviations for formulas
of \(\Lang\), not new symbols of the language. The same abbreviation
\(\exists!\), with a fresh comparison variable, is used in the statement
of the replacement schema in \cref{sec:zf-zfc-01}.
\index{uniqueness quantifier}\index{exactly one}
\end{hbdefinition}

\begin{hbtheorem}{thm:p01-w02-p01-022-uniqueness}{proved=annex:p01-w02-p01-022-uniqueness}{Counting characterization of uniqueness}
Let \(\varphi\) be a formula, \(x\) a variable, \(y\) a comparison variable
for \((\varphi,x)\), \(\mathcal M\) a structure, \(s\) an assignment, and
\(S=\varphi^{\mathcal M,s}_x\).
\begin{enumerate}
  \item \(\mathcal M,s\entails\exists^{\leq1}x\,\varphi\) iff \(S\) has at
        most one element.
  \item \(\mathcal M,s\entails\exists!x\,\varphi\) iff \(S\) has exactly one
        element.
  \item \(\exists!x\,\varphi\) is logically equivalent to
        \(\exists x\,\varphi\land\exists^{\leq1}x\,\varphi\) and to
        \(\exists x\,\forall y\,(\varphi[y/x]\leftrightarrow y=x)\).
\end{enumerate}
Since \(S\) does not depend on \(y\), the truth conditions in clauses 1 and 2
are the same for every comparison variable.
\end{hbtheorem}

\begin{hblemma}{thm:p01-w02-p01-022-freshness}{proved=annex:p01-w02-p01-022-freshness}{The comparison variable must be fresh}
Let \(x\ne y\) and let \(\varphi\) be the formula \(x=y\), in which \(y\)
occurs free. Put
\[
  \chi\defeq\exists x\,\bigl(x=y\land\forall y\,(y=y\to y=x)\bigr),
\]
the result of using \(y\) in the pattern of
\cref{def:p01-w02-p01-022-uniqueness}, since \((x=y)[y/x]\) is \(y=y\). In
every structure \(\mathcal M\) with at least two elements and every
assignment \(s\), \(\varphi^{\mathcal M,s}_x=\{s(y)\}\) has exactly one
element, while \(\mathcal M,s\not\entails\chi\).
\end{hblemma}

\begin{hbtheorem}{thm:p01-w02-p01-022-exactly-one}{proved=annex:p01-w02-p01-022-exactly-one}{Exactly one P is Q, and the P is Q}
Let \(P,Q\) be unary relation symbols and \(y\ne x\). Put
\[
\begin{aligned}
  E_1&\defeq\exists x\,\bigl(P(x)\land Q(x)\land
        \forall y\,((P(y)\land Q(y))\to y=x)\bigr),\\
  E_2&\defeq\exists x\,\bigl(P(x)\land\forall y\,(P(y)\to y=x)\land
        Q(x)\bigr).
\end{aligned}
\]
Then \(E_1\) is logically equivalent to \(\exists!x\,(P(x)\land Q(x))\),
\(E_2\entails E_1\),
\(E_1\not\entails E_2\), and \(\exists x\,(P(x)\land Q(x))\not\entails E_1\).
\end{hbtheorem}

\(E_1\) says that exactly one \(P\) is a \(Q\); \(E_2\) says that there is
exactly one \(P\) and that it is a \(Q\). ``Exactly one referee read the
draft'' is of the first form; ``the referee read the draft'' is of the
second. The uniqueness clause attaches to the noun phrase that ``exactly
one'' or ``the'' governs, and the translator records which phrase that is.

\begin{hbexample}{def:p01-w02-p01-022-referees}{Uniqueness inside a universal}
With \(R(x)\) referee, \(D(y)\) paper and \(H(x,y)\) ``\(x\) read \(y\)'',
``every referee read exactly one paper'' is
\[
  \forall x\,\bigl(R(x)\to\exists!y\,(D(y)\land H(x,y))\bigr),
\]
where the comparison variable is a third variable \(z\) and the
abbreviation expands to
\(\exists y\,(D(y)\land H(x,y)\land\forall z\,((D(z)\land H(x,z))\to
z=y))\). The paper may vary with the referee. ``One paper was read by every
referee'', in its reading ``exactly one paper was read by every referee'',
is \(\exists!y\,(D(y)\land\forall x\,(R(x)\to H(x,y)))\): the uniqueness now
concerns papers read by all referees, and the two sentences differ in
scope exactly as in the inherited example ``Quantifier order'' of
\cref{sec:levels-05}. The same two sentence shapes, with reviewers in place
of referees, are posed for translation in the inherited exercises on
first-order deduction
\cite[ch.~``First-Order Deduction'', exercise
\texttt{ex:fol-deduction-06}]{FrozenFormalLogicFormalizing}.
\end{hbexample}

\subsection{Definite descriptions and their scope}
\label{sec:p01-w02-p01-022-descriptions}

The terms over a signature (\cref{sec:fol-syntax-semantics-01}) are
variables and function applications only, so ``the \(\varphi\)'' is not a
term of \(\Lang\). It is
translated by eliminating it from the sentence in which it occurs.

\begin{hbdefinition}{def:p01-w02-p01-022-description}{Contextual elimination of a description}
Let \(\varphi\) and \(\chi\) be formulas, \(x\) a variable, and \(y\) a
comparison variable for \((\varphi,x)\). The sentence frame ``the
\(x\) such that \(\varphi\) is such that \(\chi\)'' is translated by
\[
  \mathsf{the}_x(\varphi,\chi)\defeq
  \exists x\,\bigl(\varphi\land\forall y\,(\varphi[y/x]\to y=x)\land\chi\bigr).
\]
When the frame occurs inside a larger context, such as a negation, the
elimination may be applied to the whole context or only to the part inside
it. Relative to a negation, the \emph{wide} reading of ``the \(\varphi\) is
not \(\chi\)'' is \(\mathsf{the}_x(\varphi,\neg\chi)\) and the \emph{narrow}
reading is \(\neg\,\mathsf{the}_x(\varphi,\chi)\). This is Russell's
analysis: a statement about ``the term having the property \(F\)'' asserts
that exactly one term has \(F\) and that this term has the further property
\cite[481--482, 490]{Russell1905}; his \emph{primary} and \emph{secondary}
occurrences of a denoting phrase are the wide and narrow readings
\cite[489--490]{Russell1905}. Here \(\mathsf{the}_x\) is a metalinguistic
abbreviation, not a term former of \(\Lang\).
\index{definite description!contextual elimination}
\index{definite description!scope}
\end{hbdefinition}

\begin{hbtheorem}{thm:p01-w02-p01-022-description-scope}{proved=annex:p01-w02-p01-022-description-scope}{Truth conditions and scope of eliminated descriptions}
Let \(\varphi,\chi\), \(x\) and \(y\) be as in
\cref{def:p01-w02-p01-022-description}, and let \(S=\varphi^{\mathcal M,s}_x\).
\begin{enumerate}
  \item \(\mathcal M,s\entails\mathsf{the}_x(\varphi,\chi)\) iff
        \(S=\{a\}\) for some \(a\) with
        \(\mathcal M,s[x\mapsto a]\entails\chi\).
  \item If \(S\) is empty or has at least two elements, then
        \(\mathcal M,s\not\entails\mathsf{the}_x(\varphi,\chi)\) for every
        formula \(\chi\).
  \item \(\mathsf{the}_x(\varphi,\neg\chi)\entails
        \neg\,\mathsf{the}_x(\varphi,\chi)\).
  \item \(\exists!x\,\varphi,\ \neg\,\mathsf{the}_x(\varphi,\chi)\entails
        \mathsf{the}_x(\varphi,\neg\chi)\).
  \item If \(S=\varnothing\), then
        \(\mathcal M,s\entails\neg\,\mathsf{the}_x(\varphi,\chi)\) and
        \(\mathcal M,s\not\entails\mathsf{the}_x(\varphi,\neg\chi)\).
\end{enumerate}
\end{hbtheorem}

Clause 2 is the feature Russell emphasizes: if the uniqueness condition
fails, every statement of the form ``the \(\varphi\) is \(\chi\)'' is false
\cite[482]{Russell1905}. Clauses 3 to 5 show that the scope of the
description matters exactly when \(\exists!x\,\varphi\) may fail: the wide
reading entails the narrow one, the converse holds given exactly one
\(\varphi\), and when there is no \(\varphi\) the narrow reading is true and
the wide one false. A translation record therefore states the scope of every
description that occurs under negation or another operator, and states
whether \(\exists!x\,\varphi\) belongs to the premises or to \(\Pi\).

\begin{hbwarning}{def:p01-w02-p01-022-description-warning}{A description is not a name}
Replacing ``the \(\varphi\)'' by a new constant \(c\) asserts nothing about
existence or uniqueness: \(\neg\chi(c)\) is satisfiable whether or not
anything satisfies \(\varphi\). The replacement is sound only if
\(\exists!x\,\varphi\) and \(\varphi[c/x]\) are added as premises or to
\(\Pi\) and recorded. Languages with a primitive description term former
\(\iota x\,\varphi\) must fix what an improper description denotes, and the
Russellian elimination is then one policy among several; see
\cref{sec:p02-free-logic,sec:p08-conditionals-deontic-free}.
\end{hbwarning}

\subsection{Worked example}
\label{sec:p01-w02-p01-022-worked}

\begin{hbexample}{def:p01-w02-p01-022-worked}{A translation record}
Argument: every referee who read a paper graded it; the paper by Ana was
read by a referee; therefore some referee graded the paper by Ana.

\emph{Domain and signature.} People and papers. Unary \(R\) (referee) and
\(D\) (paper); binary \(H\) (read), \(G\) (graded) and \(B\) (``\(y\) is
by \(z\)''); constant \(a\) (Ana). \(\Pi=\varnothing\): no background
sentence is needed for the verdict.

\emph{Ambiguity points.} (i) Pronoun: ``it'' is bound by ``a paper'' inside
the restrictor; recorded as a universal over referee-paper pairs
(\cref{def:p01-w02-p01-022-pronoun}). (ii) Description: ``the paper by Ana''
is eliminated with uniqueness, with wide scope in both sentences in which it
occurs; nothing in the argument places it under an operator, so no scope
alternative arises. (iii) Number: ``a referee'' and ``some referee'' assert
at least one.

\emph{Translation.} With comparison variable \(z\) for the description,
\[
\begin{aligned}
  \mathrm{P}_1&\defeq\forall x\,\forall y\,\bigl((R(x)\land D(y)\land
       H(x,y))\to G(x,y)\bigr),\\
  \mathrm{P}_2&\defeq\mathsf{the}_y\bigl(D(y)\land B(y,a),\
       \exists x\,(R(x)\land H(x,y))\bigr),\\
  \mathrm{C}&\defeq\mathsf{the}_y\bigl(D(y)\land B(y,a),\
       \exists x\,(R(x)\land G(x,y))\bigr).
\end{aligned}
\]
\Cref{thm:p01-w02-p01-022-worked} shows that \(\mathrm{P}_1,\mathrm{P}_2
\entails\mathrm{C}\). It also shows that dropping the uniqueness clause from
the second premise, which turns it into
\[
  \mathrm{P}_2'\defeq\exists y\,\bigl(D(y)\land B(y,a)\land
  \exists x\,(R(x)\land H(x,y))\bigr),
\]
breaks the argument: the conclusion asserts uniqueness, and
nothing else supplies it.
\end{hbexample}

\begin{hbtheorem}{thm:p01-w02-p01-022-worked}{proved=annex:p01-w02-p01-022-worked}{Verdicts for the worked example}
With the notation of \cref{def:p01-w02-p01-022-worked},
\(\mathrm{P}_1,\mathrm{P}_2\entails\mathrm{C}\) and
\(\mathrm{P}_1,\mathrm{P}_2'\not\entails\mathrm{C}\).
\end{hbtheorem}

\subsection{Limitations and relationships}
\label{sec:p01-w02-p01-022-limits}

The tests above decide whether two readings differ; they do not decide which
reading a speaker meant. Resolution rests on stated evidence, and a record
that cannot resolve a point reports validity per reading. Every reading here
is a first-order formula over a declared signature; a reading that the
signature cannot express, such as the collective reading of
\cref{def:p01-w02-p01-022-plurality} before collections are declared, forces
a change of signature, which is itself recorded. Tense and modality, which
step 3 of the inherited method also lists, require temporal or modal
languages and are not translated by the patterns of this section.

The section extends three inherited places: step 3 and the example
``Quantifier order'' of \cref{sec:levels-05}, the abbreviation of restricted
quantifiers in \cref{sec:fol-semantics-03}, and the propositional
translations of \cref{sec:prop-syntax-semantics-03}, where ``exactly one of
\(P\) and \(Q\)'' is the exclusive disjunction that
\cref{def:p01-w02-p01-022-uniqueness} generalizes to a formula with a free
variable. Many-sorted relativization is in \cref{sec:p02-many-sorted};
primitive description terms and their free-logic policies are in
\cref{sec:p02-free-logic,sec:p08-conditionals-deontic-free}.
